% Blueprint content for "Primes of the Form x^2 + ny^2" (Cox, 2nd ed.).
%
% Each node carries \label, \lean (the matching Lean declaration), \uses (its
% dependencies in the DAG), and \leanok on every statement whose Lean signature
% typechecks. Proof environments deliberately carry no \leanok: every Lean proof
% is currently `sorry` (this is a scaffold).
%
% See ROADMAP.md for the Mathlib audit of each node.

\chapter{Part I — From Fermat to Gauss}

This part is largely formalizable with current Mathlib (quadratic forms, \texttt{ZMod},
quadratic reciprocity). It develops the classical theory of binary quadratic forms,
their class group, and genus theory, culminating in congruence criteria for
representation by the principal form.

\section{Fermat, Euler, and quadratic reciprocity}

This section follows Cox \S1. The three Fermat theorems are obtained from Euler's
two-step strategy (Descent Step plus Reciprocity Step); the Reciprocity Step is
solved in general by quadratic reciprocity, packaged through the character
$\chi$ of Lemma 1.14 and Corollary 1.19.

\begin{theorem}[Two-square identity; Cox (1.3)]
  \label{eq:mul_sq_add_sq}
  \lean{PrimesXNY2.Fermat.mul_sq_add_sq}
  \leanok
  $(x^2+y^2)(z^2+w^2) = (xz - yw)^2 + (xw + yz)^2$.
\end{theorem}

\begin{theorem}[Identity for $x^2+ny^2$; Cox (1.6)]
  \label{eq:mul_sq_add_nsq}
  \lean{PrimesXNY2.Fermat.mul_sq_add_nsq}
  \leanok
  $(x^2+ny^2)(z^2+nw^2) = (xz - nyw)^2 + n(xw + yz)^2$.
\end{theorem}

\begin{theorem}[Cox Lemma 1.4]
  \label{lem:descent_lemma}
  \lean{PrimesXNY2.Fermat.descent_lemma}
  \leanok
  If $N$ is a sum of two relatively prime squares and a prime $q = x^2+y^2$
  divides $N$, then $N/q$ is again a sum of two relatively prime squares.
\end{theorem}
\begin{proof}
  \uses{eq:mul_sq_add_sq}
\end{proof}

\begin{theorem}[Descent Step, $n=1$; Cox \S1]
  \label{thm:descent}
  \lean{PrimesXNY2.Fermat.descent_step}
  \leanok
  If an odd prime $p$ divides $x^2+y^2$ with $\gcd(x,y)=1$, then $p$ is itself a
  sum of two squares.
\end{theorem}
\begin{proof}
  \uses{lem:descent_lemma}
  Repeated application of Lemma 1.4 (\ref{lem:descent_lemma}) descends to $p$.
\end{proof}

\begin{theorem}[Two squares; Cox Thm 1.2]
  \label{thm:fermat_two_squares}
  \lean{PrimesXNY2.Fermat.prime_sq_add_sq}
  \leanok
  An odd prime $p$ is a sum of two squares, $p = x^2 + y^2$, if and only if
  $p \equiv 1 \pmod 4$.
\end{theorem}
\begin{proof}
  \uses{thm:descent, thm:legendre_supplementary}
  Reciprocity shows $-1$ is a quadratic residue mod $p$ exactly when
  $p \equiv 1 \pmod 4$; descent (\ref{thm:descent}) then produces the
  representation.
\end{proof}

\begin{theorem}[$x^2 + 2y^2$]
  \label{thm:fermat_x2_2y2}
  \lean{PrimesXNY2.Fermat.prime_sq_add_two_sq}
  \leanok
  An odd prime $p$ satisfies $p = x^2 + 2y^2$ iff $p \equiv 1, 3 \pmod 8$.
\end{theorem}
\begin{proof}
  \uses{thm:descent, thm:legendre_supplementary}
\end{proof}

\begin{theorem}[$x^2 + 3y^2$]
  \label{thm:fermat_x2_3y2}
  \lean{PrimesXNY2.Fermat.prime_sq_add_three_sq}
  \leanok
  A prime $p > 3$ satisfies $p = x^2 + 3y^2$ iff $p \equiv 1 \pmod 3$.
\end{theorem}
\begin{proof}
  \uses{thm:descent, thm:quadratic_reciprocity}
\end{proof}

\begin{definition}[Legendre symbol]
  \label{def:legendre_symbol}
  \lean{legendreSym}
  \leanok
  For an odd prime $p$ and $a \in \mathbb Z$, the Legendre symbol $(a/p)$ is
  $0, 1, -1$ according as $p \mid a$, $a$ is a nonzero square, or a nonsquare mod
  $p$. (Provided by Mathlib.)
\end{definition}

\begin{theorem}[Cox Lemma 1.7]
  \label{lem:dvd_iff_isSquare}
  \lean{PrimesXNY2.Fermat.dvd_sq_add_nsq_iff_isSquare}
  \leanok
  For nonzero $n$ and an odd prime $p \nmid n$, $p$ divides a primitive value
  $x^2+ny^2$ iff $(-n/p) = 1$.
\end{theorem}
\begin{proof}
  \uses{def:legendre_symbol}
\end{proof}

\begin{theorem}[Euler's form of reciprocity; Cox Conjecture 1.9]
  \label{conj:euler_reciprocity}
  \lean{PrimesXNY2.Fermat.euler_reciprocity}
  \leanok
  For distinct odd primes $p,q$, $(q/p)=1$ iff $p \equiv \pm\beta^2 \pmod{4q}$ for
  some odd $\beta$.
\end{theorem}
\begin{proof}
  \uses{def:legendre_symbol}
\end{proof}

\begin{theorem}[Quadratic Reciprocity; Cox Proposition 1.10]
  \label{thm:quadratic_reciprocity}
  \lean{PrimesXNY2.Fermat.quadratic_reciprocity}
  \leanok
  For distinct odd primes $p,q$,
  $(p/q)(q/p) = (-1)^{\frac{p-1}{2}\cdot\frac{q-1}{2}}$.
\end{theorem}
\begin{proof}
  \uses{conj:euler_reciprocity, def:legendre_symbol}
\end{proof}

\begin{theorem}[Supplementary laws]
  \label{thm:legendre_supplementary}
  \lean{PrimesXNY2.Fermat.legendreSym_supplementary}
  \leanok
  $(-1/p) = (-1)^{(p-1)/2}$ and $(2/p) = (-1)^{(p^2-1)/8}$.
\end{theorem}
\begin{proof}
  \uses{def:legendre_symbol}
\end{proof}

\begin{theorem}[Cox Lemma 1.14]
  \label{lem:quadraticChar}
  \lean{PrimesXNY2.Fermat.exists_quadraticChar}
  \leanok
  For nonzero $D \equiv 0,1 \pmod 4$ there is a homomorphism
  $\chi : (\mathbb Z/D\mathbb Z)^\times \to \{\pm 1\}$ with $\chi([m]) = (D/m)$
  (Jacobi symbol) for odd $m$ prime to $D$.
\end{theorem}
\begin{proof}
  \uses{thm:quadratic_reciprocity, def:legendre_symbol}
\end{proof}

\begin{theorem}[Reciprocity Step solved; Cox Corollary 1.19]
  \label{thm:reciprocity_step}
  \lean{PrimesXNY2.Fermat.reciprocity_step}
  \leanok
  For each $n>0$ there is a set of residues mod $4n$ characterizing the odd
  primes $p \nmid n$ that divide a primitive $x^2+ny^2$.
\end{theorem}
\begin{proof}
  \uses{lem:dvd_iff_isSquare, lem:quadraticChar}
\end{proof}

\subsection{Exercises}

\begin{theorem}[Exercise 1.1(a)]
  \label{exr:s1_1_a}
  \lean{PrimesXNY2.PartI.S1.ex_1_1_a}
  \leanok
  Prove the identity (1.3).
\end{theorem}
\begin{proof}\uses{eq:mul_sq_add_sq}\end{proof}

\begin{theorem}[Exercise 1.1(b)]
  \label{exr:s1_1_b}
  \lean{PrimesXNY2.PartI.S1.ex_1_1_b}
  \leanok
  Euler's generalization
  $(ax^2+cy^2)(az^2+cw^2) = (axz - cyw)^2 + ac(xw + yz)^2$.
\end{theorem}
\begin{proof}\uses{eq:mul_sq_add_nsq}\end{proof}

\begin{theorem}[Exercise 1.2(a)]
  \label{exr:s1_2_a}
  \lean{PrimesXNY2.PartI.S1.ex_1_2_a}
  \leanok
  For $k \ge 1$, $\Delta^k$ applied to $g$ is an integral linear combination of
  $g(x), \dots, g(x+k)$.
\end{theorem}

\begin{theorem}[Exercise 1.2(b)]
  \label{exr:s1_2_b}
  \lean{PrimesXNY2.PartI.S1.ex_1_2_b}
  \leanok
  For a monic $f$ of degree $d$, $\Delta^d f = d!$.
\end{theorem}

\begin{theorem}[Exercise 1.2(c)]
  \label{exr:s1_2_c}
  \lean{PrimesXNY2.PartI.S1.ex_1_2_c}
  \leanok
  A monic integer polynomial of degree $< p$ is not identically $0 \bmod p$.
\end{theorem}
\begin{proof}\uses{exr:s1_2_a, exr:s1_2_b}\end{proof}

\begin{theorem}[Exercise 1.3(a)]
  \label{exr:s1_3_a}
  \lean{PrimesXNY2.PartI.S1.ex_1_3_a}
  \leanok
  Lemma 1.4 for the form $x^2+ny^2$.
\end{theorem}
\begin{proof}\uses{lem:descent_lemma}\end{proof}

\begin{theorem}[Exercise 1.3(b)]
  \label{exr:s1_3_b}
  \lean{PrimesXNY2.PartI.S1.ex_1_3_b}
  \leanok
  The descent works for $n=3$, $q=4$.
\end{theorem}
\begin{proof}\uses{exr:s1_3_a}\end{proof}

\begin{theorem}[Exercise 1.4(a)]
  \label{exr:s1_4_a}
  \lean{PrimesXNY2.PartI.S1.ex_1_4_a}
  \leanok
  Descent Step for $x^2+2y^2$.
\end{theorem}
\begin{proof}\uses{exr:s1_3_a}\end{proof}

\begin{theorem}[Exercise 1.4(b)]
  \label{exr:s1_4_b}
  \lean{PrimesXNY2.PartI.S1.ex_1_4_b}
  \leanok
  Descent Step for $x^2+3y^2$ (odd $p$).
\end{theorem}
\begin{proof}\uses{exr:s1_3_b}\end{proof}

\begin{theorem}[Exercise 1.5]
  \label{exr:s1_5}
  \lean{PrimesXNY2.PartI.S1.ex_1_5}
  \leanok
  If $p = 3k+1$ is prime then $(-3/p) = 1$.
\end{theorem}

\begin{theorem}[Exercise 1.6]
  \label{exr:s1_6}
  \lean{PrimesXNY2.PartI.S1.ex_1_6}
  \leanok
  Prove Lemma 1.7.
\end{theorem}
\begin{proof}\uses{lem:dvd_iff_isSquare}\end{proof}

\begin{theorem}[Exercise 1.7]
  \label{exr:s1_7}
  \lean{PrimesXNY2.PartI.S1.ex_1_7}
  \leanok
  Quadratic reciprocity in the form (1.12): $(p^*/q) = (q/p)$.
\end{theorem}
\begin{proof}\uses{thm:quadratic_reciprocity}\end{proof}

\begin{theorem}[Exercise 1.8]
  \label{exr:s1_8}
  \lean{PrimesXNY2.PartI.S1.ex_1_8}
  \leanok
  The reciprocity statement (1.13).
\end{theorem}
\begin{proof}\uses{conj:euler_reciprocity}\end{proof}

\begin{theorem}[Exercise 1.9(a)]
  \label{exr:s1_9_a}
  \lean{PrimesXNY2.PartI.S1.ex_1_9_a}
  \leanok
  For $p>3$: $(-3/p)=1 \iff p \equiv 1 \pmod 3$.
\end{theorem}
\begin{proof}\uses{lem:dvd_iff_isSquare}\end{proof}

\begin{theorem}[Exercise 1.9(b)]
  \label{exr:s1_9_b}
  \lean{PrimesXNY2.PartI.S1.ex_1_9_b}
  \leanok
  The cases $x^2+y^2$, $x^2+2y^2$ are the supplementary laws.
\end{theorem}
\begin{proof}\uses{thm:legendre_supplementary}\end{proof}

\begin{theorem}[Exercise 1.10(a)]
  \label{exr:s1_10_a}
  \lean{PrimesXNY2.PartI.S1.ex_1_10_a}
  \leanok
  $(M/m) = (N/m)$ when $M \equiv N \pmod m$.
\end{theorem}

\begin{theorem}[Exercise 1.10(b)]
  \label{exr:s1_10_b}
  \lean{PrimesXNY2.PartI.S1.ex_1_10_b}
  \leanok
  Multiplicativity of the Jacobi symbol (1.15).
\end{theorem}

\begin{theorem}[Exercise 1.10(c)]
  \label{exr:s1_10_c}
  \lean{PrimesXNY2.PartI.S1.ex_1_10_c}
  \leanok
  The supplementary laws (1.16) for the Jacobi symbol.
\end{theorem}

\begin{theorem}[Exercise 1.10(d)]
  \label{exr:s1_10_d}
  \lean{PrimesXNY2.PartI.S1.ex_1_10_d}
  \leanok
  If $M$ is a quadratic residue mod $m$ then $(M/m)=1$ (converse false).
\end{theorem}

\begin{theorem}[Exercise 1.11]
  \label{exr:s1_11}
  \lean{PrimesXNY2.PartI.S1.ex_1_11}
  \leanok
  Completion of (1.17): $(D/m) = (D/n)$ for odd $m \equiv n \pmod D$.
\end{theorem}
\begin{proof}\uses{exr:s1_10_b, exr:s1_10_c}\end{proof}

\begin{theorem}[Exercise 1.12(a)]
  \label{exr:s1_12_a}
  \lean{PrimesXNY2.PartI.S1.ex_1_12_a}
  \leanok
  $\chi([m]) = (D/m)$ is a well-defined homomorphism.
\end{theorem}
\begin{proof}\uses{lem:quadraticChar}\end{proof}

\begin{theorem}[Exercise 1.12(b)]
  \label{exr:s1_12_b}
  \notready
  $\chi([-1]) = 1$ if $D > 0$ and $-1$ if $D < 0$. \emph{UNVERIFIED: a faithful
  self-contained statement must fix the specific $\chi$ of Lemma 1.14; flagged.}
\end{theorem}

\begin{theorem}[Exercise 1.12(c)]
  \label{exr:s1_12_c}
  \notready
  Value of $\chi([2])$ when $D \equiv 1 \pmod 4$. \emph{UNVERIFIED: same reason as
  1.12(b); flagged.}
\end{theorem}

\begin{theorem}[Exercise 1.13(a)]
  \label{exr:s1_13_a}
  \lean{PrimesXNY2.PartI.S1.ex_1_13_a}
  \leanok
  Quadratic reciprocity, assuming Lemma 1.14.
\end{theorem}
\begin{proof}\uses{lem:quadraticChar}\end{proof}

\begin{theorem}[Exercise 1.13(b)]
  \label{exr:s1_13_b}
  \lean{PrimesXNY2.PartI.S1.ex_1_13_b}
  \leanok
  The supplementary laws, assuming Lemma 1.14.
\end{theorem}
\begin{proof}\uses{lem:quadraticChar}\end{proof}

\begin{theorem}[Exercise 1.14]
  \label{exr:s1_14}
  \lean{PrimesXNY2.PartI.S1.ex_1_14}
  \leanok
  When $n \equiv 3 \pmod 4$ the Reciprocity Step congruence can be taken mod $n$.
\end{theorem}
\begin{proof}\uses{thm:reciprocity_step}\end{proof}

\begin{theorem}[Exercise 1.15]
  \label{exr:s1_15}
  \lean{PrimesXNY2.PartI.S1.ex_1_15}
  \leanok
  The residue classes in $(\mathbb Z/84\mathbb Z)^\times$ with $(-21/p)=1$.
\end{theorem}
\begin{proof}\uses{thm:quadratic_reciprocity}\end{proof}

\begin{theorem}[Exercise 1.16]
  \label{exr:s1_16}
  \notready
  When $D = 4q$ ($q$ an odd prime), properties (i) and (ii) suffice to determine
  $K = \ker\chi$. \emph{UNVERIFIED: a faithful statement of the uniqueness
  characterization of the index-two subgroup is delicate; flagged.}
\end{theorem}

\section{Binary quadratic forms and reduction}

\begin{definition}[Binary quadratic form]
  \label{def:binary_qf}
  \lean{PrimesXNY2.Forms.BinaryQF}
  \leanok
  An integral binary quadratic form $f(x,y) = a x^2 + b x y + c y^2$, recorded by
  its coefficient triple $(a,b,c)$.
\end{definition}

\begin{definition}[Discriminant]
  \label{def:discriminant}
  \lean{PrimesXNY2.Forms.BinaryQF.discr}
  \leanok
  \uses{def:binary_qf}
  The discriminant of $f = (a,b,c)$ is $D = b^2 - 4ac$.
\end{definition}

\begin{definition}[Proper equivalence]
  \label{def:proper_equiv}
  \lean{PrimesXNY2.Forms.ProperlyEquivalent}
  \leanok
  \uses{def:binary_qf}
  Two forms are properly equivalent if related by a determinant-one integer change
  of variables, i.e.\ by an element of $\mathrm{SL}_2(\mathbb Z)$.
\end{definition}

\begin{definition}[Reduced form]
  \label{def:reduced}
  \lean{PrimesXNY2.Forms.BinaryQF.Reduced}
  \leanok
  \uses{def:binary_qf, def:discriminant}
  A positive definite form $(a,b,c)$ is reduced when $|b| \le a \le c$, with
  $b \ge 0$ whenever $|b| = a$ or $a = c$.
\end{definition}

\begin{theorem}[Reduction theorem; Cox Thm 2.8]
  \label{thm:reduction}
  \lean{PrimesXNY2.Forms.exists_unique_reduced}
  \leanok
  Every positive definite form is properly equivalent to a unique reduced form.
\end{theorem}
\begin{proof}
  \uses{def:reduced, def:proper_equiv}
\end{proof}

\begin{theorem}[Finiteness of the class number; Cox Thm 2.13]
  \label{thm:finite_class_number}
  \lean{PrimesXNY2.Forms.finite_reduced_of_discr}
  \leanok
  For each $D < 0$ there are only finitely many reduced forms of discriminant $D$.
\end{theorem}
\begin{proof}
  \uses{thm:reduction}
\end{proof}

\begin{definition}[Form evaluation]
  \label{def:form_eval}
  \lean{PrimesXNY2.Forms.BinaryQF.eval}
  \leanok
  \uses{def:binary_qf}
  $f(x,y) = a x^2 + b x y + c y^2$ as a function of $(x,y)$.
\end{definition}

\begin{definition}[Primitive form]
  \label{def:primitive}
  \lean{PrimesXNY2.Forms.BinaryQF.Primitive}
  \leanok
  \uses{def:binary_qf}
  A form is primitive when $\gcd(a,b,c)=1$.
\end{definition}

\begin{definition}[Positive definite]
  \label{def:posdef}
  \lean{PrimesXNY2.Forms.BinaryQF.PosDef}
  \leanok
  \uses{def:discriminant}
  $a > 0$ and $D < 0$.
\end{definition}

\begin{definition}[Indefinite]
  \label{def:indefinite}
  \lean{PrimesXNY2.Forms.BinaryQF.Indefinite}
  \leanok
  \uses{def:discriminant}
  $D > 0$.
\end{definition}

\begin{definition}[Action of an integer matrix]
  \label{def:action}
  \lean{PrimesXNY2.Forms.action}
  \leanok
  \uses{def:binary_qf}
  The change of variables $(x,y) \mapsto (px+qy, rx+sy)$ by a matrix $M$.
\end{definition}

\begin{definition}[Full equivalence; Cox (2.2)]
  \label{def:equivalent}
  \lean{PrimesXNY2.Forms.Equivalent}
  \leanok
  \uses{def:action}
  Forms related by a change of variables of determinant $\pm 1$
  ($\mathrm{GL}_2(\mathbb Z)$).
\end{definition}

\begin{theorem}[Proper equivalence is an equivalence relation]
  \label{thm:proper_equiv_equivalence}
  \lean{PrimesXNY2.Forms.properlyEquivalent_equivalence}
  \leanok
  \uses{def:proper_equiv}
\end{theorem}

\begin{theorem}[Proper equivalence preserves the discriminant]
  \label{thm:discr_invariant}
  \lean{PrimesXNY2.Forms.discr_eq_of_properlyEquivalent}
  \leanok
  \uses{def:proper_equiv, def:discriminant}
\end{theorem}

\begin{theorem}[Cox (2.4)]
  \label{eq:four_mul_eval}
  \lean{PrimesXNY2.Forms.four_mul_eval}
  \leanok
  \uses{def:form_eval, def:discriminant}
  $4a\,f(x,y) = (2ax+by)^2 - D y^2$.
\end{theorem}

\begin{definition}[Principal form]
  \label{def:principal_form}
  \lean{PrimesXNY2.Forms.principalForm}
  \leanok
  \uses{def:discriminant}
  $x^2 - (D/4)y^2$ if $D \equiv 0$, else $x^2 + xy + ((1-D)/4)y^2$.
\end{definition}

\begin{theorem}[Cox Lemma 2.3]
  \label{lem:repr_iff_equiv}
  \lean{PrimesXNY2.Genus.properlyRepresents_iff_properlyEquivalent}
  \leanok
  \uses{def:represents, def:proper_equiv}
  $f$ properly represents $m$ iff $f \sim [m,b,c]$ for some $b,c$.
\end{theorem}

\begin{theorem}[Cox Lemma 2.5]
  \label{lem:repr_iff_isSquare_general}
  \lean{PrimesXNY2.Genus.properlyRepresents_iff_isSquare_general}
  \leanok
  \uses{lem:repr_iff_equiv, def:primitive}
  For $D \equiv 0,1 \pmod 4$ and odd $m$ prime to $D$, $m$ is properly represented
  by a primitive form of discriminant $D$ iff $D$ is a square mod $m$.
\end{theorem}

\begin{theorem}[Cox Corollary 2.6]
  \label{cor:c2_6}
  \lean{PrimesXNY2.Genus.cor_2_6}
  \leanok
  \uses{lem:repr_iff_isSquare_general}
  For an odd prime $p \nmid n$, $(-n/p)=1$ iff $p$ is represented by a primitive
  form of discriminant $-4n$.
\end{theorem}

\begin{theorem}[Cox Proposition 2.15]
  \label{prop:p2_15}
  \lean{PrimesXNY2.Genus.prop_2_15}
  \leanok
  \uses{cor:c2_6, thm:reduction}
  For an odd prime $p \nmid n$, $(-n/p)=1$ iff $p$ is represented by one of the
  reduced forms of discriminant $-4n$.
\end{theorem}

\begin{theorem}[Cox Theorem 2.16]
  \label{thm:t2_16}
  \lean{PrimesXNY2.Genus.thm_2_16}
  \leanok
  \uses{prop:p2_15, lem:quadraticChar}
  For negative $D \equiv 0,1 \pmod 4$ and odd prime $p \nmid D$, $(D/p)=1$
  (i.e.\ $[p] \in \ker\chi$) iff $p$ is represented by a reduced form of disc $D$.
\end{theorem}

\begin{theorem}[Composition identity (2.30)]
  \label{eq:comp_2_30}
  \lean{PrimesXNY2.Forms.comp_identity_2_30}
  \leanok
\end{theorem}

\begin{theorem}[Composition identity (2.31)]
  \label{eq:comp_2_31}
  \lean{PrimesXNY2.Forms.comp_identity_2_31}
  \leanok
\end{theorem}

\begin{theorem}[Landau; Cox Theorem 2.18]
  \label{thm:landau}
  \lean{PrimesXNY2.Forms.class_number_one}
  \leanok
  \uses{def:reduced, def:primitive}
  $h(-4n)=1$ iff $n \in \{1,2,3,4,7\}$.
\end{theorem}

\section{The form class group and Dirichlet composition}

\begin{definition}[Form class group]
  \label{def:form_class_group}
  \lean{PrimesXNY2.Forms.FormClassGroup}
  \leanok
  \uses{def:proper_equiv, def:discriminant}
  $C(D)$ is the set of proper equivalence classes of forms of discriminant $D$.
\end{definition}

\begin{definition}[Dirichlet composition]
  \label{def:dirichlet_composition}
  \lean{PrimesXNY2.Forms.compose}
  \leanok
  \uses{def:form_class_group}
  The group law on $C(D)$ given by Dirichlet composition of forms.
\end{definition}

\begin{theorem}[$C(D)$ is a finite abelian group; Cox Thm 3.9]
  \label{thm:form_class_group_is_group}
  \lean{PrimesXNY2.Forms.isCommGroup}
  \leanok
  Under Dirichlet composition, $C(D)$ is a finite abelian group with identity the
  principal class.
\end{theorem}
\begin{proof}
  \uses{def:dirichlet_composition, thm:reduction}
\end{proof}

\section{Genus theory and representation}

\begin{definition}[Proper representation]
  \label{def:represents}
  \lean{PrimesXNY2.Genus.ProperlyRepresents}
  \leanok
  \uses{def:binary_qf}
  $f$ properly represents $m$ if $m = f(x,y)$ with $\gcd(x,y) = 1$.
\end{definition}

\begin{definition}[Genus]
  \label{def:genus}
  \lean{PrimesXNY2.Genus.genus}
  \leanok
  \uses{def:form_class_group}
  The genus of a form: the residues mod $D$ it represents, constant on each genus.
\end{definition}

\begin{definition}[Principal genus]
  \label{def:principal_genus}
  \lean{PrimesXNY2.Genus.principalGenus}
  \leanok
  \uses{def:genus}
  The genus of the principal form, equivalently the subgroup of squares in $C(D)$.
\end{definition}

\begin{theorem}[Cox Lemma 2.5, prime case]
  \label{thm:representation_residue}
  \lean{PrimesXNY2.Genus.properlyRepresents_iff_isSquare}
  \leanok
  \uses{def:represents}
  A prime $p \nmid D$ is properly represented by some form of discriminant $D$ iff
  $D$ is a quadratic residue mod $p$. (Prime case of Lemma 2.5,
  \ref{lem:repr_iff_isSquare_general}.)
\end{theorem}

\begin{theorem}[Cox Lemma 2.24]
  \label{lem:l2_24}
  \lean{PrimesXNY2.Genus.principalForm_values_subgroup}
  \leanok
  \uses{def:principal_form, def:genus}
  For negative $D \equiv 0,1 \pmod 4$, the residues in
  $(\mathbb Z/D\mathbb Z)^\times$ represented by the principal form form a
  subgroup $H$.
\end{theorem}

\begin{theorem}[Cox Lemma 2.25]
  \label{lem:l2_25}
  \lean{PrimesXNY2.Genus.properlyRepresents_coprime}
  \leanok
  \uses{def:form_eval}
  Every form properly represents some value relatively prime to a given $M$.
\end{theorem}

\begin{theorem}[Cox Theorem 2.26]
  \label{thm:t2_26}
  \lean{PrimesXNY2.Genus.thm_2_26}
  \leanok
  \uses{lem:l2_24, def:genus, thm:reduction}
  An odd prime $p \nmid D$ lies in the genus of $g$ iff $p$ is represented by a
  reduced form of disc $D$ in the same genus.
\end{theorem}

\begin{theorem}[Cox Corollary 2.27]
  \label{thm:genus_congruence}
  \lean{PrimesXNY2.Genus.represents_principal_iff_congruence}
  \leanok
  \uses{def:principal_genus, lem:l2_24}
  For $n>0$ and odd prime $p \nmid n$, $p$ is represented by a form of disc $-4n$
  in the principal genus iff $p \equiv \beta^2$ or $\beta^2 + n \pmod{4n}$.
  (Replaces an earlier vacuous/false statement.)
\end{theorem}

\subsection{Exercises}

\begin{theorem}[Exercise 2.1]
  \label{exr:s2_1}
  \lean{PrimesXNY2.PartI.S2.ex_2_1}
  \leanok
  Any represented value is $d^2$ times a properly represented value.
\end{theorem}
\begin{proof}\uses{def:represents}\end{proof}

\begin{theorem}[Exercise 2.2(a)]
  \label{exr:s2_2_a}
  \lean{PrimesXNY2.PartI.S2.ex_2_2_a}
  \leanok
  Equivalence and proper equivalence are equivalence relations.
\end{theorem}
\begin{proof}\uses{def:equivalent, thm:proper_equiv_equivalence}\end{proof}

\begin{theorem}[Exercise 2.2(b)]
  \label{exr:s2_2_b}
  \lean{PrimesXNY2.PartI.S2.ex_2_2_b}
  \leanok
  Improper equivalence is not an equivalence relation.
\end{theorem}
\begin{proof}\uses{def:action}\end{proof}

\begin{theorem}[Exercise 2.2(c)]
  \label{exr:s2_2_c}
  \lean{PrimesXNY2.PartI.S2.ex_2_2_c}
  \leanok
  Equivalent forms represent the same numbers.
\end{theorem}
\begin{proof}\uses{def:equivalent}\end{proof}

\begin{theorem}[Exercise 2.2(d)]
  \label{exr:s2_2_d}
  \lean{PrimesXNY2.PartI.S2.ex_2_2_d}
  \leanok
  A form equivalent to a primitive form is primitive.
\end{theorem}
\begin{proof}\uses{def:equivalent, def:primitive}\end{proof}

\begin{theorem}[Exercise 2.3]
  \label{exr:s2_3}
  \lean{PrimesXNY2.PartI.S2.ex_2_3}
  \leanok
  $D_f = (\det M)^2 D_g$ under a change of variables.
\end{theorem}
\begin{proof}\uses{def:action, def:discriminant}\end{proof}

\begin{theorem}[Exercise 2.4(a)]
  \label{exr:s2_4_a}
  \lean{PrimesXNY2.PartI.S2.ex_2_4_a}
  \leanok
  A form of positive discriminant represents both signs.
\end{theorem}
\begin{proof}\uses{eq:four_mul_eval}\end{proof}

\begin{theorem}[Exercise 2.4(b)]
  \label{exr:s2_4_b}
  \lean{PrimesXNY2.PartI.S2.ex_2_4_b}
  \leanok
  A form of negative discriminant is definite (sign of $a$).
\end{theorem}
\begin{proof}\uses{eq:four_mul_eval}\end{proof}

\begin{theorem}[Exercise 2.5]
  \label{exr:s2_5}
  \lean{PrimesXNY2.PartI.S2.ex_2_5}
  \leanok
  Corollary 2.6 for arbitrary discriminant.
\end{theorem}
\begin{proof}\uses{cor:c2_6}\end{proof}

\begin{theorem}[Exercise 2.6]
  \label{exr:s2_6}
  \lean{PrimesXNY2.PartI.S2.ex_2_6}
  \leanok
  A reduced form properly equivalent to $126x^2+74xy+13y^2$ exists.
\end{theorem}
\begin{proof}\uses{thm:reduction}\end{proof}

\begin{theorem}[Exercise 2.7]
  \label{exr:s2_7}
  \lean{PrimesXNY2.PartI.S2.ex_2_7}
  \leanok
  The bound (2.9).
\end{theorem}
\begin{proof}\uses{def:reduced}\end{proof}

\begin{theorem}[Exercise 2.8(a)]
  \label{exr:s2_8_a}
  \lean{PrimesXNY2.PartI.S2.ex_2_8_a}
  \leanok
  The minima (2.11) of a reduced form.
\end{theorem}
\begin{proof}\uses{def:reduced}\end{proof}

\begin{theorem}[Exercise 2.8(b)]
  \label{exr:s2_8_b}
  \lean{PrimesXNY2.PartI.S2.ex_2_8_b}
  \leanok
  Uniqueness part of Theorem 2.8 (exceptional cases).
\end{theorem}
\begin{proof}\uses{thm:reduction}\end{proof}

\begin{theorem}[Exercise 2.9(a)]
  \label{exr:s2_9_a}
  \notready
  Verify the table (2.14) of class numbers and reduced forms. \emph{UNVERIFIED:
  computational verification of a table, not a single proposition; flagged.}
\end{theorem}

\begin{theorem}[Exercise 2.9(b)]
  \label{exr:s2_9_b}
  \notready
  Compute the reduced forms of discriminants $-3,-15,-24,-31,-52$. \emph{UNVERIFIED:
  a per-discriminant enumeration, not a single faithfully-stateable proposition;
  flagged.}
\end{theorem}

\begin{theorem}[Exercise 2.10(a)]
  \label{exr:s2_10_a}
  \lean{PrimesXNY2.PartI.S2.ex_2_10_a}
  \leanok
  Reduction for indefinite forms.
\end{theorem}
\begin{proof}\uses{def:indefinite, thm:reduction}\end{proof}

\begin{theorem}[Exercise 2.10(b)]
  \label{exr:s2_10_b}
  \lean{PrimesXNY2.PartI.S2.ex_2_10_b}
  \leanok
  Bound $4a^2 \le D$.
\end{theorem}
\begin{proof}\uses{exr:s2_10_a}\end{proof}

\begin{theorem}[Exercise 2.10(c)]
  \label{exr:s2_10_c}
  \lean{PrimesXNY2.PartI.S2.ex_2_10_c}
  \leanok
  Finiteness of the indefinite class number.
\end{theorem}
\begin{proof}\uses{exr:s2_10_b}\end{proof}

\begin{theorem}[Exercise 2.11]
  \label{exr:s2_11}
  \lean{PrimesXNY2.PartI.S2.ex_2_11}
  \leanok
  The result (2.17) for $x^2+7y^2$.
\end{theorem}
\begin{proof}\uses{thm:t2_16, thm:quadratic_reciprocity}\end{proof}

\begin{theorem}[Exercise 2.12(a)]
  \label{exr:s2_12_a}
  \lean{PrimesXNY2.PartI.S2.ex_2_12_a}
  \leanok
  Non-prime-power factorization.
\end{theorem}

\begin{theorem}[Exercise 2.12(b)]
  \label{exr:s2_12_b}
  \lean{PrimesXNY2.PartI.S2.ex_2_12_b}
  \leanok
  $h(-32)=2$ and $h(-124)=3$.
\end{theorem}
\begin{proof}\uses{def:reduced}\end{proof}

\begin{theorem}[Exercise 2.13]
  \label{exr:s2_13}
  \lean{PrimesXNY2.PartI.S2.ex_2_13}
  \leanok
  The result (2.19) for $x^2+5y^2$.
\end{theorem}
\begin{proof}\uses{thm:t2_16, thm:quadratic_reciprocity}\end{proof}

\begin{theorem}[Exercise 2.14]
  \label{exr:s2_14}
  \lean{PrimesXNY2.PartI.S2.ex_2_14}
  \leanok
  The residues (2.20) of $x^2+5y^2$ and $2x^2+2xy+3y^2$.
\end{theorem}
\begin{proof}\uses{lem:l2_24}\end{proof}

\begin{theorem}[Exercise 2.15]
  \label{exr:s2_15}
  \lean{PrimesXNY2.PartI.S2.ex_2_15}
  \leanok
  The result (2.23) for $x^2+14y^2$, $2x^2+7y^2$.
\end{theorem}
\begin{proof}\uses{thm:t2_26}\end{proof}

\begin{theorem}[Exercise 2.16]
  \label{exr:s2_16}
  \lean{PrimesXNY2.PartI.S2.ex_2_16}
  \leanok
  The principal form for $D \equiv 1 \pmod 4$ is reduced of disc $D$.
\end{theorem}
\begin{proof}\uses{def:principal_form, def:reduced}\end{proof}

\begin{theorem}[Exercise 2.17(a)]
  \label{exr:s2_17_a}
  \lean{PrimesXNY2.PartI.S2.ex_2_17_a}
  \leanok
  Even values force $D \equiv 1 \pmod 8$.
\end{theorem}
\begin{proof}\uses{lem:repr_iff_equiv}\end{proof}

\begin{theorem}[Exercise 2.17(b)]
  \label{exr:s2_17_b}
  \lean{PrimesXNY2.PartI.S2.ex_2_17_b}
  \leanok
  Represented units lie in $\ker\chi$.
\end{theorem}
\begin{proof}\uses{lem:repr_iff_isSquare_general}\end{proof}

\begin{theorem}[Exercise 2.17(c)]
  \label{exr:s2_17_c}
  \lean{PrimesXNY2.PartI.S2.ex_2_17_c}
  \leanok
  $H$ is the subgroup of squares (principal form values).
\end{theorem}
\begin{proof}\uses{def:principal_form, lem:l2_24}\end{proof}

\begin{theorem}[Exercise 2.17(d)]
  \label{exr:s2_17_d}
  \notready
  The values represented by $f$ form a coset of $H$. \emph{UNVERIFIED: a faithful
  coset statement requires fixing the square subgroup and a representative;
  flagged (same delicacy as Lemma 2.24(ii)).}
\end{theorem}

\begin{theorem}[Exercise 2.18(a)]
  \label{exr:s2_18_a}
  \lean{PrimesXNY2.PartI.S2.ex_2_18_a}
  \leanok
  Some of $f(1,0),f(0,1),f(1,1)$ is prime to $p$.
\end{theorem}
\begin{proof}\uses{def:primitive}\end{proof}

\begin{theorem}[Exercise 2.18(b)]
  \label{exr:s2_18_b}
  \lean{PrimesXNY2.PartI.S2.ex_2_18_b}
  \leanok
  Proof of Lemma 2.25.
\end{theorem}
\begin{proof}\uses{lem:l2_25}\end{proof}

\begin{theorem}[Exercise 2.19]
  \label{exr:s2_19}
  \notready
  Work out the genus theory of Theorem 2.26 for $D=-15,-24,-31,-52$.
  \emph{UNVERIFIED: produces four separate congruence theorems, not one
  proposition; flagged.}
\end{theorem}

\begin{theorem}[Exercise 2.20]
  \label{exr:s2_20}
  \notready
  Corollary 2.27 for negative $D \equiv 1 \pmod 4$. \emph{UNVERIFIED: Cox does not
  give the explicit congruence for this case and I will not guess it; flagged.}
\end{theorem}

\begin{theorem}[Exercise 2.21]
  \label{exr:s2_21}
  \lean{PrimesXNY2.PartI.S2.ex_2_21}
  \leanok
  The first theorem of (2.28): $x^2+6y^2$.
\end{theorem}
\begin{proof}\uses{thm:genus_congruence}\end{proof}

\begin{theorem}[Exercise 2.22]
  \label{exr:s2_22}
  \lean{PrimesXNY2.PartI.S2.ex_2_22}
  \leanok
  The composition identity (2.31).
\end{theorem}
\begin{proof}\uses{eq:comp_2_31}\end{proof}

\begin{theorem}[Exercise 2.23(a)]
  \label{exr:s2_23_a}
  \lean{PrimesXNY2.PartI.S2.ex_2_23_a}
  \leanok
  $p \equiv 1 \pmod 8 \Rightarrow (-2/p)=1$.
\end{theorem}

\begin{theorem}[Exercise 2.23(b)]
  \label{exr:s2_23_b}
  \lean{PrimesXNY2.PartI.S2.ex_2_23_b}
  \leanok
  $p \equiv 3 \pmod 8$: $(2/p)=1$ and representation by a disc-$8$ form.
\end{theorem}
\begin{proof}\uses{cor:c2_6}\end{proof}

\begin{theorem}[Exercise 2.23(c)]
  \label{exr:s2_23_c}
  \lean{PrimesXNY2.PartI.S2.ex_2_23_c}
  \leanok
  Forms of discriminant $8$ are $\pm(x^2-2y^2)$.
\end{theorem}
\begin{proof}\uses{thm:reduction}\end{proof}

\begin{theorem}[Exercise 2.23(d)]
  \label{exr:s2_23_d}
  \lean{PrimesXNY2.PartI.S2.ex_2_23_d}
  \leanok
  $p = \pm(x^2-2y^2) \Rightarrow p \equiv \pm 1 \pmod 8$.
\end{theorem}

\begin{theorem}[Exercise 2.24]
  \label{exr:s2_24}
  \lean{PrimesXNY2.PartI.S2.ex_2_24}
  \leanok
  Legendre's theorem on $ax^2+by^2+cz^2=0$.
\end{theorem}

\begin{theorem}[Exercise 2.24(a)]
  \label{exr:s2_24_a}
  \lean{PrimesXNY2.PartI.S2.ex_2_24_a}
  \leanok
  $x^2+py^2-qz^2=0$ has a nontrivial solution.
\end{theorem}
\begin{proof}\uses{exr:s2_24}\end{proof}

\begin{theorem}[Exercise 2.24(b)]
  \label{exr:s2_24_b}
  \lean{PrimesXNY2.PartI.S2.ex_2_24_b}
  \leanok
  The mod-$4$ obstruction (contradiction proving reciprocity).
\end{theorem}

\begin{theorem}[Exercise 2.25]
  \label{exr:s2_25}
  \lean{PrimesXNY2.PartI.S2.ex_2_25}
  \leanok
  Forms are properly equivalent iff their opposites are.
\end{theorem}
\begin{proof}\uses{def:proper_equiv}\end{proof}

\begin{theorem}[Exercise 2.26]
  \label{exr:s2_26}
  \lean{PrimesXNY2.PartI.S2.ex_2_26}
  \leanok
  The four compositions lie in distinct classes.
\end{theorem}
\begin{proof}\uses{def:proper_equiv, exr:s2_6}\end{proof}

\begin{theorem}[Exercise 2.27(a)]
  \label{exr:s2_27_a}
  \lean{PrimesXNY2.PartI.S2.ex_2_27_a}
  \leanok
  A prime represented by two forms of the same disc: they are equivalent.
\end{theorem}
\begin{proof}\uses{lem:repr_iff_equiv, def:equivalent}\end{proof}

\begin{theorem}[Exercise 2.27(b)]
  \label{exr:s2_27_b}
  \lean{PrimesXNY2.PartI.S2.ex_2_27_b}
  \leanok
  A reduced form equivalent to $x^2+ny^2$ equals it.
\end{theorem}
\begin{proof}\uses{thm:reduction}\end{proof}

\chapter{Part II — Class Field Theory}

This part states the structural bridge between forms and ideals, then the deep
class field theory inputs (Artin reciprocity, Čebotarev) that current Mathlib does
not yet provide; those are stated to be cited or assumed. See ROADMAP.md.

\section{Orders in imaginary quadratic fields}

\begin{definition}[Order]
  \label{def:quad_order}
  \lean{PrimesXNY2.Order.QuadOrder}
  \leanok
  An order $\mathcal O = \mathbb Z + f\,\mathcal O_K$ of conductor $f$ in an
  imaginary quadratic field $K$.
\end{definition}

\begin{definition}[Discriminant of an order]
  \label{def:order_discriminant}
  \lean{PrimesXNY2.Order.QuadOrder.discr}
  \leanok
  \uses{def:quad_order}
  The discriminant $D = f^2 d_K$ of the order $\mathcal O$.
\end{definition}

\begin{definition}[Ideal class group of an order]
  \label{def:ideal_class_group_order}
  \lean{PrimesXNY2.Order.QuadOrder.idealClassGroup}
  \leanok
  \uses{def:quad_order}
  $C(\mathcal O)$: proper fractional $\mathcal O$-ideals modulo principal ideals.
\end{definition}

\begin{theorem}[$C(\mathcal O)$ is a finite abelian group; Cox Prop 7.19]
  \label{thm:ideal_class_group_finite}
  \lean{PrimesXNY2.Order.QuadOrder.idealClassGroup_isCommGroup}
  \leanok
  \uses{def:ideal_class_group_order}
  The ideal class group of an order is a finite abelian group.
\end{theorem}

\section{The bridge: forms and ideals}

\begin{theorem}[Bridge; Cox Thm 7.7]
  \label{thm:form_ideal_bridge}
  \lean{PrimesXNY2.Bridge.formClassGroup_equiv_idealClassGroup}
  \leanok
  For the order $\mathcal O$ of discriminant $D < 0$, the form class group $C(D)$ is
  isomorphic, as a group, to the ideal class group $C(\mathcal O)$.
\end{theorem}
\begin{proof}
  \uses{def:form_class_group, def:ideal_class_group_order, def:order_discriminant}
\end{proof}

\begin{theorem}[Bridge respects composition]
  \label{thm:bridge_composition}
  \lean{PrimesXNY2.Bridge.bridge_respects_composition}
  \leanok
  The bridge intertwines Dirichlet composition with ideal multiplication.
\end{theorem}
\begin{proof}
  \uses{thm:form_ideal_bridge, def:dirichlet_composition}
\end{proof}

\section{Class fields, Artin reciprocity, Čebotarev (deep)}

\begin{definition}[Hilbert class field]
  \label{def:hilbert_class_field}
  \lean{PrimesXNY2.RingClassField.hilbertClassField}
  \leanok
  \uses{def:quad_order}
  The maximal unramified abelian extension $H$ of $K$.
\end{definition}

\begin{definition}[Ring class field]
  \label{def:ring_class_field}
  \lean{PrimesXNY2.RingClassField.ringClassField}
  \leanok
  \uses{def:quad_order}
  The abelian extension $L/K$ with $\mathrm{Gal}(L/K) \cong C(\mathcal O)$.
\end{definition}

\begin{theorem}[Artin reciprocity; Cox Thm 8.2 — deep]
  \label{thm:artin_reciprocity}
  \lean{PrimesXNY2.RingClassField.artinReciprocity}
  \leanok
  The Artin map induces an isomorphism $C(\mathcal O) \cong \mathrm{Gal}(L/K)$.
  \emph{Deep: cite or assume.}
\end{theorem}
\begin{proof}
  \uses{def:ring_class_field, def:ideal_class_group_order}
\end{proof}

\begin{definition}[Complete splitting]
  \label{def:splits_completely}
  \lean{PrimesXNY2.RingClassField.SplitsCompletely}
  \leanok
  \uses{def:ring_class_field}
  A prime $p$ splits completely in the ring class field of $\mathcal O$.
\end{definition}

\begin{theorem}[Čebotarev density; deep]
  \label{thm:chebotarev}
  \lean{PrimesXNY2.RingClassField.chebotarev_splitting}
  \leanok
  The primes splitting completely in the ring class field are infinite (density
  $1/|C(\mathcal O)|$). \emph{Deep: cite or assume.}
\end{theorem}
\begin{proof}
  \uses{def:splits_completely}
\end{proof}

\section{The Main Theorem}

\begin{definition}[Class polynomial]
  \label{def:class_polynomial}
  \lean{PrimesXNY2.MainTheorem.classPolynomial}
  \leanok
  \uses{def:ring_class_field}
  The integer polynomial $f_n(x)$ whose roots generate the ring class field of
  $\mathbb Z[\sqrt{-n}]$.
\end{definition}

\begin{theorem}[Main Theorem, splitting form; Cox Thm 9.2]
  \label{thm:main_splitting}
  \lean{PrimesXNY2.MainTheorem.prime_repr_iff_splits}
  \leanok
  A prime $p \nmid n$ is of the form $x^2 + n y^2$ iff $p$ splits completely in the
  ring class field.
\end{theorem}
\begin{proof}
  \uses{thm:form_ideal_bridge, thm:artin_reciprocity, def:splits_completely, thm:chebotarev}
\end{proof}

\begin{theorem}[Main Theorem, class polynomial form; Cox Cor 9.4]
  \label{thm:main_class_poly}
  \lean{PrimesXNY2.MainTheorem.prime_repr_iff_classPoly_root}
  \leanok
  $p = x^2 + n y^2$ iff $(-n/p) = 1$ and $f_n(x) \equiv 0 \pmod p$ has a root.
\end{theorem}
\begin{proof}
  \uses{thm:main_splitting, def:class_polynomial}
\end{proof}

\chapter{Part III — Complex Multiplication}

The deepest part: elliptic functions, the $j$-invariant, modular functions for
$\mathrm{SL}_2(\mathbb Z)$, complex multiplication, and Weber functions with
explicit class polynomials and Shimura reciprocity (2nd ed.). Almost entirely
to be formalized; see ROADMAP.md.

\section{Elliptic functions and the j-invariant}

\begin{definition}[Lattice]
  \label{def:lattice}
  \lean{PrimesXNY2.Elliptic.Lattice}
  \leanok
  A lattice $L = \mathbb Z\omega_1 + \mathbb Z\omega_2 \subseteq \mathbb C$.
\end{definition}

\begin{definition}[Weierstrass $\wp$]
  \label{def:weierstrass_p}
  \lean{PrimesXNY2.Elliptic.weierstrassP}
  \leanok
  \uses{def:lattice}
  The Weierstrass $\wp$-function of a lattice $L$.
\end{definition}

\begin{definition}[$j$-invariant of a lattice]
  \label{def:j_invariant_lattice}
  \lean{PrimesXNY2.Elliptic.jInvariant}
  \leanok
  \uses{def:lattice}
  $j(L) = 1728\, g_2^3 / (g_2^3 - 27 g_3^2)$.
\end{definition}

\begin{theorem}[Weierstrass differential equation]
  \label{thm:weierstrass_ode}
  \lean{PrimesXNY2.Elliptic.weierstrassP_differential_eq}
  \leanok
  $(\wp')^2 = 4\wp^3 - g_2 \wp - g_3$.
\end{theorem}
\begin{proof}
  \uses{def:weierstrass_p}
\end{proof}

\section{Modular functions and complex multiplication}

\begin{definition}[The $j$-function]
  \label{def:j_function}
  \lean{PrimesXNY2.Modular.jFunction}
  \leanok
  The Klein $j$-function $j(\tau)$ on the upper half-plane.
\end{definition}

\begin{definition}[Modular function]
  \label{def:modular_function}
  \lean{PrimesXNY2.Modular.IsModularFunction}
  \leanok
  A modular function for $\mathrm{SL}_2(\mathbb Z)$.
\end{definition}

\begin{theorem}[$j$ is modular]
  \label{thm:j_modular}
  \lean{PrimesXNY2.Modular.jFunction_modular}
  \leanok
  \uses{def:j_function}
  $j(\gamma\tau) = j(\tau)$ for $\gamma \in \mathrm{SL}_2(\mathbb Z)$.
\end{theorem}
\begin{proof}
  \uses{def:modular_function}
\end{proof}

\begin{definition}[CM point]
  \label{def:cm_point}
  \lean{PrimesXNY2.Modular.IsCMPoint}
  \leanok
  An imaginary quadratic $\tau$ (a complex multiplication point).
\end{definition}

\begin{theorem}[CM values are integral; Cox Thm 11.1]
  \label{thm:cm_integral}
  \lean{PrimesXNY2.Modular.isIntegral_jFunction_of_cm}
  \leanok
  If $\tau$ is a CM point then $j(\tau)$ is an algebraic integer.
\end{theorem}
\begin{proof}
  \uses{def:cm_point, def:j_function}
\end{proof}

\begin{theorem}[$j$ generates the ring class field; Cox Thm 11.1]
  \label{thm:j_generates}
  \lean{PrimesXNY2.Modular.jFunction_generates_ringClassField}
  \leanok
  There is a CM point whose $j$-value is an algebraic integer generating the ring
  class field $L = K(j(\tau))$.
\end{theorem}
\begin{proof}
  \uses{thm:cm_integral, def:ring_class_field}
\end{proof}

\section{Weber functions and explicit class polynomials}

\begin{definition}[Weber function]
  \label{def:weber_function}
  \lean{PrimesXNY2.Weber.weber}
  \leanok
  A Weber modular function $\mathfrak f(\tau)$.
\end{definition}

\begin{definition}[Weber class polynomial]
  \label{def:weber_class_polynomial}
  \lean{PrimesXNY2.Weber.weberClassPolynomial}
  \leanok
  \uses{def:weber_function}
  The reduced class polynomial built from Weber values.
\end{definition}

\begin{theorem}[Weber matches Hilbert class polynomial mod $p$]
  \label{thm:weber_root}
  \lean{PrimesXNY2.Weber.weberClassPolynomial_root_iff}
  \leanok
  The Weber class polynomial has a root mod $p$ iff the Hilbert class polynomial
  does, so it may replace $f_n$ in the criterion.
\end{theorem}
\begin{proof}
  \uses{def:weber_class_polynomial, def:class_polynomial}
\end{proof}

\begin{theorem}[Shimura reciprocity; 2nd ed.]
  \label{thm:shimura}
  \lean{PrimesXNY2.Weber.shimuraReciprocity}
  \leanok
  The idele class group acts on CM values of modular functions, computing the
  Galois action on Weber values. \emph{Deep: cite or assume.}
\end{theorem}
\begin{proof}
  \uses{def:cm_point, thm:artin_reciprocity}
\end{proof}
